% Standalone example: compile with xelatex or lualatex from this directory,
% or copy this directory into your project and use \usetheme[path=<dir>/]{ucl2025}.
\documentclass[aspectratio=169]{beamer}
\usetheme{ucl2025}
\renewcommand{\uclfooter}{Footer text}

\title{Main headline in 42pt UCL Sans SemiBold, no more than 2--3 lines}
\author{Name of presenter}
\institute{Faculty, department or unit}
\date{\today}

\begin{document}

\begin{frame}
  \titlepage
\end{frame}

\begin{frame}{Title in 36pt UCL Sans SemiBold, no more than 1--2 lines}
  Use UCL Sans Regular no smaller than 24pt for body text, so it stays
  readable and accessible.

  Keep paragraphs short.
  \begin{itemize}
    \item using bullet points
    \item often makes text
    \item easier to read
  \end{itemize}
\end{frame}

\begin{frame}[dark]{Title in 36pt UCL Sans SemiBold, no more than 1--2 lines}
  The dark mode version of a content slide.
\end{frame}

\section{Divider}
\begin{frame}[divider]{Divider slide title in 36pt UCL Sans SemiBold}
  This slide can be used to introduce a new section.
\end{frame}

\begin{frame}[stripe]{Title in 36pt UCL Sans SemiBold, no more than 1--2 lines}
  \begin{columns}[T]
    \column{0.35\textwidth}
      \subheading{Subheading in 30pt UCL Sans Regular}
    \column{0.6\textwidth}
      Use UCL Sans Regular no smaller than 24pt for body text.
  \end{columns}
\end{frame}

\begin{frame}[divider]{Thank you}
  your.name@ucl.ac.uk
\end{frame}

\end{document}
